\chapter{System Architecture}

\section{System Context}

The Mini Payment Gateway is built around one central gateway core service. That
service coordinates merchant transaction requests, provider outcome evidence,
merchant notifications, and internal operational review. Two separate
dashboards sit around the core: one for internal operations and one for
merchant-facing visibility.

\begin{figure}[!htbp]
\centering
\resizebox{\textwidth}{!}{%
\begin{tikzpicture}[node distance=12mm and 18mm]
  \node[actor] (merchantBackend) {Merchant backend};
  \node[box, below=of merchantBackend, text width=32mm] (checkout) {Customer checkout\\QR display};
  \node[actor, below=of checkout] (customer) {Customer banking\\app};
  \node[actor, below=of customer] (provider) {Provider callback\\source};
  \node[actor, below=6mm of provider] (trigger) {Scheduled\\trigger};

  \node[box, right=34mm of provider, text width=38mm] (gateway) {Gateway core service};
  \node[box, left=of gateway, text width=32mm, yshift=-31mm] (opsDashboard) {Internal operations dashboard};
  \node[box, below=of opsDashboard, text width=32mm] (merchantDashboard) {Merchant dashboard};
  \node[actor, left=of opsDashboard, yshift=-8mm] (opsUser) {Admin / Ops user};
  \node[actor, left=of merchantDashboard] (merchantUser) {Merchant dashboard\\user};
  \node[db, right=of gateway] (store) {Persistent\\data store};
  \node[box, below=of store, text width=34mm] (webhookTarget) {Merchant callback\\endpoint};

  \draw[arrow] (merchantBackend.east) -- (gateway.west);
  \draw[arrow] (gateway.west) to[out=150,in=25] (merchantBackend.east);
  \draw[arrow] (merchantBackend.south) -- (checkout.north);
  \draw[arrow] (customer.north) -- (checkout.south);
  \draw[arrow] (provider.east) -- (gateway.west);
  \draw[arrow] (trigger.east) -- (gateway.west);
  \draw[arrow] (opsUser.east) -- (opsDashboard.west);
  \draw[arrow] (merchantUser.east) -- (merchantDashboard.west);
  \draw[arrow] (opsDashboard.east) -- (gateway.west);
  \draw[arrow] (merchantDashboard.east) -- (gateway.west);
  \draw[arrow] (gateway.east) -- (store.west);
  \draw[arrow] (gateway.east) to[out=10,in=180] (webhookTarget.west);
\end{tikzpicture}
}
\caption{Logical system context. The gateway core service sits between merchant
requests, provider evidence, scheduled processing, dashboard access, and
merchant notification delivery.}
\end{figure}

\newpage
\section{Logical Components}

\begin{table}[!htbp]
\centering
\caption{Logical components and responsibilities}
\begin{tabularx}{\textwidth}{L{0.24\textwidth}L{0.22\textwidth}Y}
\toprule
\textbf{Component} & \textbf{Role} & \textbf{Primary responsibility} \\
\midrule
Gateway core service & Central application service & Handles payment and refund requests, processes external outcome evidence, creates merchant notifications, and exposes internal and merchant-facing data views. \\
\bodyrowrule
Persistent data store & Durable data layer & Stores merchant profiles, transaction state, callback evidence, notification history, reconciliation items, audit history, and user-account information. \\
\bodyrowrule
Internal operations dashboard & Internal user interface & Supports merchant review, credential management, transaction investigation, reconciliation handling, and operational recovery workflows. \\
\bodyrowrule
Merchant dashboard & Merchant user interface & Supports merchant-scoped visibility into payments, refunds, notifications, profile information, credentials, and analytics. \\
\bodyrowrule
Merchant backend & External integration client & Sends payment and refund requests on behalf of merchant business workflows. \\
\bodyrowrule
Demo merchant application & Local reference integration & Keeps the merchant secret in server memory, signs Merchant API requests, serves a customer checkout, simulates provider callbacks, verifies gateway webhooks, and exposes the resulting order state. \\
\bodyrowrule
Customer checkout / QR display & Merchant-owned checkout surface & Displays the QR payload or image returned by the gateway through the merchant backend. \\
\bodyrowrule
Customer banking app & External payment channel & Scans the VietQR image displayed by the merchant checkout and initiates the real-world bank transfer outside the gateway UI. \\
\bodyrowrule
Provider callback source & External evidence source & Confirms payment and refund outcomes after external processing. \\
\bodyrowrule
Scheduled trigger & Time-based orchestrator & Activates expiration handling and due notification retry processing. \\
\bottomrule
\end{tabularx}
\end{table}

\newpage
\section{Authentication, Authorization, And Trust Boundaries}

\begin{table}[!htbp]
\centering
\caption{Trust boundaries and access models}
\begin{tabularx}{\textwidth}{L{0.20\textwidth}L{0.21\textwidth}L{0.24\textwidth}Y}
\toprule
\textbf{Actor} & \textbf{Access surface} & \textbf{Authentication model} & \textbf{Scope source} \\
\midrule
Merchant backend & Merchant transaction API & Signed merchant integration credentials & The authenticated merchant integration identity. \\
\bodyrowrule
Admin or Ops user & Internal operations API, including merchant dashboard account provisioning & Internal session authentication & The authenticated internal role and permissions. \\
\bodyrowrule
Admin user & Internal-user administration and high-risk merchant actions & Internal session authentication & The authenticated internal role with elevated administrative permissions. \\
\bodyrowrule
Merchant dashboard user & Merchant dashboard API & Merchant session authentication & The merchant account associated with the signed-in dashboard user. \\
\bodyrowrule
Provider callback source & Outcome callback surface & Provider id, timestamp, and HMAC callback signature & The authenticated provider identity and the referenced payment or refund evidence carried in the callback payload. \\
\bottomrule
\end{tabularx}
\end{table}

This separation is one of the system's key architectural choices. Merchant
backends authenticate as integrations, internal users authenticate for
operational control, and merchant dashboard users authenticate for
merchant-scoped visibility. After a merchant dashboard user signs in, the
system derives merchant scope from the authenticated session rather than from
user-supplied identifiers. This prevents cross-merchant access through
manipulated requests and avoids reusing integration credentials for human
dashboard login.

\newpage
\section{Key Request Flows}

Three request flows illustrate the main design decisions in the system:

\begin{enumerate}
  \item \textbf{Merchant creates a payment}: the merchant backend submits a
  signed payment request. The gateway checks merchant readiness and active
  VietQR account configuration, creates a new pending payment or reuses an
  equivalent pending request, and returns QR data for checkout.
  \item \textbf{Internal operator prepares merchant access}: an internal user
  registers the merchant, reviews readiness information, issues credentials,
  and, when needed, provisions merchant dashboard access for human merchant
  representatives.
  \item \textbf{Merchant reads analytics}: a signed-in merchant dashboard user
  requests analytics for a selected reporting window. The gateway aggregates
  merchant-owned payment, refund, and notification data and returns a scoped
  analytical view.
\end{enumerate}

\section{VietQR Pilot Flow}

The gateway remains API-only and does not host a production checkout page.
This branch includes a separate local demo merchant application that receives
QR data through the Merchant API and displays it to the customer. It represents
the behavior of a real merchant backend without moving merchant secrets into
browser JavaScript. The important system boundary is that QR generation is
controlled by gateway configuration, while the customer's bank app and the
provider callback simulator represent the external payment world.

\begin{figure}[!htbp]
\centering
\resizebox{\textwidth}{!}{%
\begin{tikzpicture}[node distance=9mm and 11mm]
  \node[actor, text width=30mm] (ops) {Admin / Ops};
  \node[box, right=of ops, text width=35mm] (config) {Configure active VietQR receiving account};
  \node[actor, below=of ops, text width=30mm] (merchant) {Merchant backend};
  \node[box, right=of merchant, text width=35mm] (payment) {Create signed payment request};
  \node[box, right=of payment, text width=38mm] (qr) {Gateway returns \code{qr_reference}, \code{qr_content}, and PNG QR};
  \node[actor, below=of qr, text width=34mm] (customer) {Customer banking app};
  \node[actor, below=of customer, text width=34mm] (provider) {Provider callback source};
  \node[box, left=of provider, text width=38mm] (callback) {Signed callback finalizes payment};
  \node[box, left=of callback, text width=35mm] (webhook) {Worker delivers merchant webhook};

  \draw[arrow] (ops) -- (config);
  \draw[arrow] (config) -- (payment);
  \draw[arrow] (merchant) -- (payment);
  \draw[arrow] (payment) -- (qr);
  \draw[arrow] (qr) -- (customer);
  \draw[arrow] (customer) -- (provider);
  \draw[arrow] (provider) -- (callback);
  \draw[arrow] (callback) -- (webhook);
\end{tikzpicture}
}
\caption{VietQR pilot flow. Internal operations configure the merchant's
receiving account, the Merchant API returns scannable QR evidence, external
payment confirmation arrives through a signed callback, and the worker handles
merchant notification delivery.}
\end{figure}

Scanning the QR only transfers the encoded bank, account, amount, and transfer
purpose into the customer's banking application. It does not call the gateway.
The provider callback informs the gateway of the external result, and the
worker's signed webhook informs the merchant. The complete executable sequence
is maintained in \code{docs/architecture/diagrams/e2e-payment-demo.puml}.
